\documentclass[11pt]{article}
\usepackage{amsmath,amssymb,bm,cancel}
\usepackage[margin=2.5cm]{geometry}
\usepackage{booktabs}

\title{The glide: the interpolation inside the RAYEN head}
\author{}
\date{}

\begin{document}
\maketitle

\section{The formula}

A gap of $N$ steps. We know $p_R$, the real value just after it. At each step:
\[
\boxed{\;P_{\text{new}} = P_{\text{old}} + \frac{p_R - P_{\text{old}}}{m}\;}
\qquad
m = \text{moves left (the $k$ gap steps left, plus the final one onto } p_R).
\]
\emph{Distance left, divided by moves left.} That is all it says.

Expand it into a cleaner form by splitting the fraction:
\[
P_{\text{new}} = P_{\text{old}} + \frac{p_R}{m} - \frac{P_{\text{old}}}{m}
= P_{\text{old}}\left(1 - \frac1m\right) + \frac{p_R}{m}
= P_{\text{old}}\,\frac{m-1}{m} + \frac{p_R}{m}. \tag{1}
\]

Write $D = p_R - P$ for the distance still to go. Subtract (1) from $p_R$:
\[
D_{\text{new}} = p_R - P_{\text{new}}
= p_R - P_{\text{old}}\frac{m-1}{m} - \frac{p_R}{m}
= p_R\,\frac{m-1}{m} - P_{\text{old}}\,\frac{m-1}{m}
= D_{\text{old}}\,\frac{m-1}{m}. \tag{2}
\]

Everything below is (1) and (2).

\section{Why the steps come out equal}

Step size is $s = D/m$. After the move, $D$ becomes $D\frac{m-1}{m}$ by (2) and
$m$ becomes $m-1$. So the next step size is
\[
s_{\text{next}} = \frac{D_{\text{new}}}{m-1}
= \frac{D\,\dfrac{m-1}{m}}{m-1}
= \frac{D}{m}
= s .
\]
The $(m-1)$ cancels. \textbf{Every move is the same size} --- a straight line ---
and nobody put the straightness in; it fell out of the cancellation.

\section{Why it always lands exactly on $p_R$}

Apply (2) once per step, with $m$ counting down from $N+1$:
\[
D_N = D_0\cdot\frac{N}{N+1}\cdot\frac{N-1}{N}\cdot\frac{N-2}{N-1}\cdots\frac{1}{2}.
\]
Every numerator cancels the next denominator:
\[
D_N = D_0\cdot\frac{\cancel{N}}{N+1}\cdot\frac{\cancel{N-1}}{\cancel{N}}\cdots\frac{1}{\cancel{2}}
= \frac{D_0}{N+1}.
\]
So after the last gap step the distance remaining is $D_0/(N+1)$ --- exactly one
more move of the standard size $s = D_0/(N+1)$. It lands on the nose.

\paragraph{Check.} $p_L=100$, $p_R=500$, $N=4$. Then $D_0=400$, $s=400/5=80$:
\[
100 \to 180 \to 260 \to 340 \to 420 \to 500 .
\]

\section{What an error does}

This is the part that matters. Suppose at some step the answer is knocked off by
$e$ --- something else in the fleet ran out of room, so this generator absorbed
the difference. Take two trajectories, one clean and one knocked, and subtract
(1) from itself:
\[
P'_{\text{new}} - P_{\text{new}}
= \left(P'_{\text{old}}\frac{m-1}{m} + \frac{p_R}{m}\right)
- \left(P_{\text{old}}\frac{m-1}{m} + \frac{p_R}{m}\right)
= (P'_{\text{old}} - P_{\text{old}})\,\frac{m-1}{m}.
\]
The $p_R/m$ terms cancel, so the error obeys the same rule (2) as the distance.
Telescoping it over $r$ steps, exactly as in Section 3:
\[
\boxed{\;e_r = e\cdot\frac{m-r}{m}\;}
\]

Read that carefully. It is \textbf{linear in $r$}: the error does not decay
quickly, it bleeds off in a straight line and only reaches zero at the seam,
where $r=m$.

\paragraph{Check.} Knock of $e=120$ with $m=4$ moves left:
\begin{center}
\begin{tabular}{ccccc}
\toprule
$r$ & clean & knocked & error & $e\frac{m-r}{m}$ \\
\midrule
1 & 260 & 350 & 90 & $120\cdot\frac34 = 90$ \\
2 & 340 & 400 & 60 & $120\cdot\frac24 = 60$ \\
3 & 420 & 450 & 30 & $120\cdot\frac14 = 30$ \\
4 & 500 & 500 & 0 & $120\cdot\frac04 = 0$ \\
\bottomrule
\end{tabular}
\end{center}

\paragraph{Two consequences.}
\begin{itemize}
\item The error is \textbf{never corrected}. It vanishes at the end only because
      $p_R$ is pinned, not because anything noticed the mistake.
\item Injecting an error at \emph{every} step and summing the linear decays
      builds a \textbf{hump} --- small at both seams, largest in the middle.
\end{itemize}

\section{The measurement}

Filling $18{:}00$--$21{:}00$ on $185$ test days:

\begin{center}
\begin{tabular}{lc}
\toprule
 & MAE (MW) \\
\midrule
straight line $p_L \to p_R$, drawn once & $55.1$ \\
the glide inside the head & $63.1$ \\
\bottomrule
\end{tabular}
\end{center}

Same formula, $8$ MW worse, because the glide is re-aimed from a knocked
position at every one of the $36$ steps. The predicted hump is visible: gas OCGT
sits flat near $168$ MW in the data, and the glide arches to $236$ MW mid-window.

\section{The ramp limit}

Moves are trimmed to what the plant can physically do:
$c = \operatorname{clip}(P_{\text{new}}, \ell, u)$. Trimming cannot cost us the
landing, because $\ell$ and $u$ carry the guard
\[
\ell \;\ge\; p_R - m\cdot(\text{fastest rise}),
\qquad
u \;\le\; p_R + m\cdot(\text{fastest fall}),
\]
i.e.\ \emph{never stand anywhere you could not get back from in the $m$ moves you
have left.}

\vspace{1em}
\noindent\rule{\linewidth}{0.4pt}\\[4pt]
\footnotesize Balance snap, ray, and the HardNet projection with proofs:
\texttt{planB/two\_heads\_explained\_full.tex}.

\end{document}
